\documentclass[11pt,a4paper]{article}
\usepackage[T1]{fontenc}
\usepackage[utf8]{inputenc}
\usepackage[margin=20mm,headheight=14pt,headsep=8mm,footskip=10mm]{geometry}
\usepackage{mathpazo}
\usepackage[scale=0.92]{tgheros}
\usepackage{microtype}
\usepackage{xcolor}
\usepackage{fancyhdr}
\usepackage{lastpage}
\definecolor{ink}{HTML}{15252C}
\definecolor{accent}{HTML}{176B70}
\definecolor{muted}{HTML}{53636B}
\setlength{\parindent}{0pt}
\setlength{\parskip}{5pt}
\linespread{1.04}
\widowpenalty=10000
\clubpenalty=10000
\brokenpenalty=10000
\pagestyle{fancy}
\fancyhf{}
\fancyhead[L]{\sffamily\footnotesize\color{muted}CLICK2FUTURE}
\fancyhead[R]{\sffamily\footnotesize\color{muted}AI6132\quad /\quad GROUP 07}
\fancyfoot[L]{\sffamily\footnotesize\color{muted}Contribution statement\quad |\quad 12 November 2026}
\fancyfoot[R]{\sffamily\footnotesize\color{muted}\thepage\ /\ \pageref{LastPage}}
\renewcommand{\headrulewidth}{0pt}
\newcommand{\member}[2]{%
  \par\vspace{10pt}{\sffamily\bfseries\fontsize{13.5}{16}\selectfont\color{accent}#1}\par
  {\sffamily\small\color{muted}#2}\par\vspace{2pt}}
\newcommand{\evidence}[1]{\par{\small\color{muted}\textit{Source anchors:} #1}\par}
\begin{document}
\color{ink}
{\sffamily\bfseries\fontsize{25}{29}\selectfont Contribution statement}\par
{\large Click2Future: action-conditioned tabletop futures}\par
{\sffamily\small\color{muted}AI6132 Generative AI for Visual Synthesis\quad |\quad Group 7}\par
\vspace{5pt}
{\color{accent}\rule{\linewidth}{0.7pt}}\par
\vspace{3pt}

This statement records the completed implementation, experiments, analysis and
writing of the four members. PRAMOD SANJAY is the group coordinator. Shared
review and integration are described separately from the original implementation;
commit counts are not used as contribution percentages.

\member{PRAMOD SANJAY}{Simulator, data generation and physical training objectives}
Implemented the seeded Pymunk tabletop simulator, renderer and force/action
contract. Built deterministic rollouts with instance masks, trajectories,
physical contact labels and replay validation, together with dataset generation,
loading and scene-group splits that prevent near-identical scenes crossing
training and evaluation. Added the controlled directional, object-collision and
wall-contact cases and their contract checks. Wrote the report's simulator and
physics-assumption subsection.

Implemented the optional contact-transition and contact-depth-change training
losses and matched control/candidate recipes. Added the reference-free
physics-guided candidate-selection study and its fair uniform-subset control;
the negative accuracy result is retained. Implemented the free-motion dynamics
residual, optional momentum correction and collision-impulse residual, with
matched configurations, tests and compact evidence. Recorded the component
comparisons and exploratory limitations rather than treating enforced geometry
or momentum constraints as learned physical reliability.

\evidence{\texttt{sim/}, dataset generation and validation; contact-dynamics
training; \texttt{physics\_selection.py}, \texttt{momentum\_block.py};
PRs 80, 93 and 94; D036, D038 and D039.}

\member{KAPOOR ANANYA}{CVAE experiments, action diagnostics and contact measurement}
Set up TC2 on a second account and added the macOS/Linux setup helper.
Independently reproduced the RGB and flow CVAE pilots, ran the
$\beta$-by-latent-size validation study, and repeated the flow CVAE with two
additional training seeds. Implemented and measured force, full-horizon motion
and changed-selection probes. Traced the flow CVAE's short-window ``stationary''
result to its slower motion, and investigated missing force scaling and
unselected-object drift alongside the training data's force and selection
coverage. Wrote the corresponding report analysis and the E0 held-out discussion.

Added a known-overlap positive control for the hull-penetration audit and a
pose-based contact-timing measure with multiple hidden-physics references.
Extended that measure to the frozen D028 projection outputs and analyzed its
sub-pixel contact standoff. Reviewed collision/evaluation changes, including
the label-only hidden-physics control, and reviewed teammate model and analysis
changes.

\evidence{CVAE study and probe scripts; \texttt{penetration\_control.py},
\texttt{pose\_contact.py}; report CVAE and held-out analysis;
PRs 54, 59, 87, 88 and 92.}

\newpage
\member{MAITI SHRIJEET}{Diffusion models, experimental integration and interactive systems}
Set up the approved TC2 workspace, scheduled GPU probes and checkpoint-continuation checks. Built pixel video diffusion training and seeded sampling, numeric and spatial
action conditioning, EMA/checkpoint continuation and inference provenance
checks. Developed the supervised object-trajectory diffusion pilot and its
rigid RGB renderer, then implemented and evaluated the E5/E6 contact,
stationarity and paired-action follow-ups. Ran and documented the shared-world,
temporal-memory, coverage, local-surface and prior-decoder comparisons, the
fixed-teacher refinement repeats and the fresh-scene D018 evaluation. Retained
failed attempts and regressions rather than selecting only successful outputs.

Implemented the interactive synthetic-image comparison: object selection and
force input, four seeded futures, synchronized playback, pixel-derived tracks,
pinned/reversed actions and separately labelled simulator references. Built
the photograph interface and allocated TC2 worker paths for the externally
pretrained Force Prompting and StreamForce models. Verified checkpoint copies,
prepared the inference weights and native runtime, and checked GPU inference,
paired actions, interruption recovery, cancellation and partial-output retention.
Preserved every declared sample and kept those external-model results separate
from the controlled-model evaluation.

Independently replayed retained poses and generated videos, checked result and
artifact hashes, and backed up native outputs. Reviewed teammate code and
measurement claims; corrected penetration-control boundary populations and
excluded empty leave-one-out contact comparisons. Consolidated the report,
verified contribution attribution and prepared reproducible source packages.
Implemented and tested the portable recorded comparison, including offline
errors and Windows runtime selection. Updated the report sources against the
merged experiment records and checked the rendered deliverables.

\evidence{Diffusion training/sampling, \texttt{object\_motion.py}, contact,
world, memory, coverage, surface and prior studies, \texttt{demo/} and photo workers;
commits \texttt{6289a3e}, \texttt{fbca399}, \texttt{df2e9d7}, \texttt{5c21e30},
\texttt{4d5c903}, \texttt{f1238eb}, \texttt{4c0735c},
\texttt{4914433}, \texttt{68ec626} and \texttt{196dae6}.}

\member{BOBAN ATHUL KRISHNA}{Baseline models, evaluation and physical-quality studies}
Implemented the CVAE baseline, generated-video evaluation, tracker calibration,
selection controls, E4 and uncertainty/contact audits, and source packaging.
Ran the E10 routing comparison and D017 held-out replication, and implemented
and evaluated the D019-D030 physical-quality studies with the physics-assisted
baseline on TC2. Developed the multi-reference energy score and D031-D035
loss-scaling, projection and training-seed studies, including the failed raw-loss
replication and the narrower projected comparison.

Retrained and checked the D038/D039 comparisons on TC2, recording checkpoint,
corpus and scorer provenance and the component-level limitations. Reviewed
teammate implementations, experiment records and report changes, and checked
their reproduction and source-package consistency.

\evidence{\texttt{train\_cvae.py}, \texttt{evaluation/},
\texttt{submission.py}, \texttt{compare\_checkpoints.py};
PRs 76, 78 and 91; D017, D019-D035 and TC2 reproduction records for D038/D039.}

\vfill
{\small\color{muted}\textit{Attribution.} The photograph extension uses externally
pretrained models; their original training and model authorship are not claimed
as team contributions. Shared review and integration do not replace the original
authors' implementation credits.}
\end{document}
